\section{Integration of Assembly Wrapping}\label{sec:wrapping}
\subsection{Existing Technique and Contribution}
Assembly wrapping is an established technique based on alternative
instruction interpretations~\cite{assembly_wrapping}. Its inclusion in
this article does not claim invention of the underlying idea. The
extension contributes an implementation in Zone, interaction with the
engine's instruction and data representations, and a separate empirical
assessment.

The relevant conceptual property is that the instruction interpretation
obtained by a particular static decoding need not be the instruction
stream followed during execution. This property makes it especially
important to identify which representation is used by a similarity or
coverage measurement. A statistic from the static view and one from
the intended execution view answer different questions.

\subsection{Integration and Supported Inputs}
Wrapping is integrated with the native rewriting infrastructure and
depends on an instruction representation adequate for the selected
input. It must also interact with the final data layout and the
references affected by changes in instruction placement. Construction
can be refused when the required input interpretation is incomplete
or an unsupported case is encountered.

Some instructions or local regions are handled differently from the
main wrapped representation. The assessment therefore records the
proportion represented by wrapping, retained regions, output size,
construction failures, and functional outcomes. A claim about the
whole output must not be inferred solely from the fraction of
instructions accepted by one internal stage.

The experiment distinguishes wrapping alone, the chosen combined
native condition, and the corresponding preparation-only controls.
Their names identify complete configurations, including data handling
and architecture. Implementation availability and successful build
are reported separately from empirical functionality.

\subsection{x86 Port and Development Status}
The x64 implementation is available for assessment. The x86 port is
under active development at the time of this draft. Register width,
instruction encodings, and addressing assumptions require
architecture-specific treatment, so an x64 outcome is not transferred
to the x86 column of a result table.

The result layout reserves separate x86 entries. Those entries remain
marked as development work until a frozen version and its assessment
are available. If the port changes during a campaign, its new results
are associated with the new version; they are not silently appended
to results from an earlier implementation.

\subsection{Assessment Boundary}
Wrapping is assessed with the same separation of construction,
functionality, and representation measurements used for other methods.
In particular, a failure to construct an output remains part of the
overall outcome. Reporting functionality only for constructed binaries
can describe accepted artifacts but does not describe availability
across all attempted cases. Section~\ref{sec:methodology} defines both
denominators.
